\documentclass[aspectratio=169]{beamer}
\usetheme{Madrid}
\setbeamercovered{transparent}
\usepackage{amsfonts}
\usepackage{pifont}
\usepackage{newpxmath}
\usepackage{newpxtext}

\date{February 5, 2025}

\begin{document}

\begin{frame}{Linux From Scratch -- Day 1}
	\small
	\begin{alertblock}{Important!}
		\vspace{0.5em}
		Sign in at \url{https://z.umn.edu/at-lfs}!
		We need to measure event participation to get funding (pizza).
		\vspace{0.5em}
	\end{alertblock}
	\vspace{0.5em}
	\begin{block}{Semesterly Goals}
		\vspace{0.5em}
		\begin{itemize}
			\item Go through Linux From Scratch (\url{https://www.linuxfromscratch.org/}), \textit{maybe} some of Beyond Linux From Scratch
			\item Learn about and discuss (userspace) parts of a Linux system that aren't talked about in classes
			\item Learn about and discuss what's actually involved in creating and maintaining a Linux distribution
		\end{itemize}
		\vspace{0.5em}
	\end{block}
\end{frame}

\begin{frame}{Organization}
	\begin{itemize}
		\item Meeting every other Wednesday night, location will be in the Discord
		      \onslide<+->{
			      \begin{itemize}
				      \item There's some trouble with room reservations, but it will \textit{probably} continue to be here.
			      \end{itemize}
		      }
		\item New event, so no food yet
		      \onslide<+->{
			      \begin{itemize}
				      \item There \alert{\textbf{will}} be pizza next LFS meeting, since there won't be a speaker next week in the UKO general meeting.
			      \end{itemize}
		      }
		\item There is a Discord channel, \texttt{\#linux-from-scratch}. Give the introduction post a wave (or introduce yourself)!
		      \onslide<+->{
			      \begin{itemize}
				      \item If you run into trouble with the steps, you should post here.
				      \item There might be a Discord role later, so we can make announcements easier.
			      \end{itemize}
		      }
		\item Homework...
		      \onslide<+->{
			      \begin{itemize}
				      \item LFS is long and there's a lot of waiting for stuff to compile -- I'm considering asking people to do some parts over the course of the week
				      \item At the end (once we've seen how much is involved in those parts), I'll ask for opinions on this -- if not, we can go at a slower pace and have no homework, but maybe not enough time to do everything
			      \end{itemize}
		      }
	\end{itemize}
\end{frame}

\begin{frame}
	\begin{block}{Today's Goals}
		\begin{itemize}
			\item Discuss what LFS is and why we care about it
			\item Make sure everyone has a machine that can work through LFS
			      \onslide<2->{
				      \begin{itemize}
					      \item LFS assumes you have a Linux machine; that should Just Work.
					      \item I \textit{expect} LFS2 on Windows to work.
					      \item I \textit{expect} UTM (or another VM) on macOS to work.
					      \item Make sure it's an x86\_64 environment! Run \texttt{uname -m} to check.
				      \end{itemize}
			      }
			\item Cover chapters 1-4, start chapter 5
		\end{itemize}
	\end{block}
\end{frame}

\begin{frame}{Chapter 1: Introduction}
	(Sorta ``flying over'' this one.)

	\vspace{1em}

	Pull up the book at \url{https://z.umn.edu/uko-lfs} while I'm doing this.
	We're doing version 12.2-systemd, if you're finding it on your own.

	\vspace{1em}

	\begin{block}{What and Why is Linux From Scratch}
		\begin{itemize}
			\item Online book that goes through the process of creating a full, working Linux system without a distro.
			\item We compile \textit{every} piece of software we use, so we can see all of them.
			\item We have to configure them and patch them to work together in a few places.
			      \vspace{1em}
			\item This gives us a great opportunity to discuss all of them!
		\end{itemize}
	\end{block}
\end{frame}

\begin{frame}{Chapter 2: Preparing the Host System}
	\begin{itemize}
		\item We don't want to have to make a partition, so we can use a \textit{loop device} instead.

		      \vspace{1em}

		      \texttt{truncate -s 32G lfs.img} \\
		      \texttt{mkfs -t ext4 lfs.img} \\
		      \texttt{sudo losetup -{}-partscan -{}-find -{}-show lfs.img} \\
		      \texttt{mkdir lfs} \\
		      \texttt{sudo mount /dev/loop0 lfs} \\
		      \texttt{export LFS="\$(pwd)/lfs"}
	\end{itemize}
\end{frame}

\end{document}
